\begin{frame}{Distributional shifts}
    Call $\mathbf X$ a fitnite set of data points and $\mathbf y$ is the corresponding set of outputs.
    $P^{\text{tr}}$ is the distribution of the train data, whereas $P^{\text{te}}$ is the distribution of the test data, both are sampled from $\mathbf X$.
    \medskip

    \textbf{Definition 3. No shifts - In–Distribution (ID)} \cite{tamang2025handlingoutofdistributiondatasurvey} \label{def3} \\
    \begin{itemize}
    \item $\mathbf P^{\text{te}}(\mathbf X) = \mathbf P^{\text{tr}}(\mathbf X)$
    \item $\mathbf P^{\text{te}}(\mathbf y \mid \mathbf X) = \mathbf P^{\text{tr}}(\mathbf y \mid \mathbf X)$
    \end{itemize}
    \medskip

    \textbf{Definition 4. Covariate shifts} \cite{MORENOTORRES2012521} \label{def4} \\
    \begin{itemize}
    \item $\mathbf P^{\text{te}}(\mathbf X) \neq \mathbf P^{\text{tr}}(\mathbf X)$
    \item $\mathbf P^{\text{te}}(\mathbf y \mid \mathbf X) = \mathbf P^{\text{tr}}(\mathbf y \mid \mathbf X)$
    \end{itemize}
    \medskip

    \textbf{Definition 5. Concept/Semantic shifts} \cite{MORENOTORRES2012521} \label{def5} \\
    \begin{itemize}
    \item $\mathbf P^{\text{te}}(\mathbf X) = \mathbf P^{\text{tr}}(\mathbf X)$
    \item $\mathbf P^{\text{te}}(\mathbf y \mid \mathbf X) \neq \mathbf P^{\text{tr}}(\mathbf y \mid \mathbf X)$
    \end{itemize}
\end{frame}